\chapter{内部数据资产与实践记录}
\chapterlead{从 9 月 22 日第一次提交到现在，仓库里有 98 次提交、114 个运行目录。这一章把它们整理出来，包括没做好的部分。}

\section{数据资产}

表~\ref{tab:pr-assets} 列出团队目前积累的数据与运行记录。所有原始文档都登记了来源、页数和 SHA-256；所有引擎运行都有输入与产物的哈希清单；
所有模型会话都保存了完整事件流。这些记录既是复现的依据，也是我们改进系统时的主要信息来源。

\begin{table}[htbp]
\centering
\caption{团队内部的数据资产与运行记录（截至 2026 年 10 月 7 日）}
\label{tab:pr-assets}
\small
\begin{tabularx}{\linewidth}{@{}l r L@{}}
\toprule
\hd{类别} & \hd{数量} & \hd{说明} \\
\midrule
真实原始文档 & 7 份 & 贵州茅台、宁德时代、招商轮船 2025 年报（已核验官方来源）；恒瑞医药、中际旭创 2025 年报，宇树科技招股书上会稿，农夫山泉港股年报（待核验来源）。共 1,685 页 \\
\rowband
样例公司 & 1 套 & 虚构的“星河智造”：年报、三季报、研报草稿、一致预期、业绩预告、可比公司，由脚本生成，内部完全自洽；35 个金标准字段、7 处埋错 \\
抽取文件 & 5 家 & 经校验的 \file{facts.verified.json}、\file{extras.json}、估值假设，随代码提交（宇树、中际在工作区内） \\
\rowband
引擎运行目录 & 114 个 & 估值 42 次、全流程 19 次、样例评测 53 次；每个目录含 \file{events.jsonl} 与 \file{manifest.json} \\
研报工作区 & 5 个 & 40 次模型会话的完整事件流与用量记录；资料登记合计 246 个来源、968 条带引文的事实；21 份深度专题笔记 \\
\rowband
交付研报 & 9 份 & 三个流程版本的研报 PDF、估值工作簿与人工复核清单 \\
盲评记录 & 27 次 & 正式盲评 7 次，校准盲评 20 次，均保存逐条判定、理由与证据 \\
\rowband
人工质检 & 33 条 & 三家 v0 研报中评审判“不通过”条目的人工复核意见，经导出页面收集 \\
评分标准 & 3 版 & v0、v1、v2 的条目文件与修订记录 \\
\bottomrule
\end{tabularx}
\end{table}

\section{交付物}

图~\ref{fig:pr-gallery} 是两份交付研报的部分页面。左上是贵州茅台研报首页，其余是中际旭创研报的若干页：首页、核心观点、业绩回顾、深度专题与估值。
每份研报都有配套的估值工作簿和人工复核清单。

\begin{figure}[htbp]
\centering
\newcommand{\thumb}[2]{\fcolorbox{rule}{white}{\includegraphics[page=#2,width=0.225\linewidth]{figures/#1}}}
\thumb{report_mt.pdf}{1}\hfill\thumb{report_zjxc.pdf}{1}\hfill\thumb{report_zjxc.pdf}{3}\hfill\thumb{report_zjxc.pdf}{9}\\[6pt]
\thumb{report_zjxc.pdf}{10}\hfill\thumb{report_zjxc.pdf}{14}\hfill\thumb{report_zjxc.pdf}{19}\hfill\thumb{report_zjxc.pdf}{23}
\caption{交付研报页面节选。左上为贵州茅台（2026 年 10 月 5 日，评级“中性”，DCF 基准每股价值 1,305.28 元）；其余为中际旭创（v2 版，盲评 90 分）。}
\label{fig:pr-gallery}
\end{figure}

\section{成本与耗时}

每个模型会话的轮数、token 和折算费用都记录在工作区里（表~\ref{tab:pr-cost}）。费用是 Claude Code 按 API 价格折算的等值，不是实际账单。

\begin{table}[htbp]
\centering
\caption{五个工作区的会话用量（页数为交付 PDF 实际页数，含目录与附录）}
\label{tab:pr-cost}
\footnotesize
\setlength{\tabcolsep}{4pt}
\begin{tabular}{@{}lrrrrrrrr@{}}
\toprule
\hd{工作区} & \hd{会话} & \hd{有效时间} & \hd{墙钟跨度} & \hd{等值费用} & \hd{输入 token} & \hd{输出 token} & \hd{页数} & \hd{图表} \\
 & & （分钟） & & （美元） & （百万） & （千） & & \\
\midrule
贵州茅台 1005 & 6 & 67.9 & 5:40 & 27.19 & 43.85 & 356.7 & 42 & 83 \\
\rowband
宁德时代 1005 & 6 & 63.9 & 6:18 & 23.59 & 36.24 & 326.5 & 42 & 102 \\
招商轮船 1005 & 7 & 64.7 & 11:10 & 27.42 & 47.20 & 322.6 & 30 & 70 \\
\rowband
中际旭创 & 10 & 144.0 & 17:01 & 42.09 & 83.79 & 424.4 & 30 & 68 \\
宇树科技 & 11 & 201.6 & 22:00 & 49.42 & 106.60 & 489.0 & 30 & 67 \\
\bottomrule
\end{tabular}
\end{table}

一次通过的 1005 批次三家，平均有效时间 65.5 分钟、费用 26.07 美元，折合每张图表 0.23～0.39 美元。三家的总时间非常接近（63.9～67.9 分钟），
说明在五会话结构下，单家公司的成本主要由流程决定，受公司差异影响不大。墙钟时间更长（5.7～11.2 小时），因为会话之间有排队和等待。
中际旭创和宇树科技贵出七到九成，是因为它们跑在流程改版期间，记录里包含旧流程的会话、一次返工和一次排版重做。

按角色看（图~\ref{fig:pr-role}），分析师和排版设计师最贵，复核编辑最便宜。这与各自的工作量相符：分析师要逐页读年报、跑校验循环、写四五份专题；
排版设计师要写几百行脚本并反复渲染；复核编辑“只删不增”，用的又是 Sonnet。

\begin{figure}[htbp]
\centering
\begin{tikzpicture}
\begin{axis}[xbar, bar width=9pt, width=0.80\linewidth, height=5.0cm, xmin=0, xmax=9,
  symbolic y coords={复核编辑,研究助理,主笔,排版设计师,分析师}, ytick=data,
  xlabel={平均等值费用（美元）}, nodes near coords, nodes near coords style={font=\scriptsize},
  axis y line*=left, axis x line*=bottom, y tick label style={font=\small\heiti}]
\addplot[fill=ink!75,draw=ink] coordinates {(2.76,复核编辑) (4.68,研究助理) (5.60,主笔) (5.83,排版设计师) (7.19,分析师)};
\end{axis}
\end{tikzpicture}
\caption{1005 批次各角色的平均费用。有效时间依次为 6.7、10.8、14.4、16.9、16.7 分钟。}
\label{fig:pr-role}
\end{figure}

计算引擎不调用模型，单次估值运行 5～11 秒；评审每次约 1.75 美元、3 分钟，是生成成本的 4\%～9\%。

\section{踩过的坑}

下面这些问题都真实发生过，有的已经解决，有的还没有。我们把它们记下来，因为它们比成功案例更能说明系统的边界。

\begin{story}[规则通道第一次跑真实年报，宁德时代一个数也没取到]
规则通道在虚构样例上一直工作得很好，第一次跑三家真实年报时，301 个模型核验值只佐证了 115 个，宁德时代是 0。原因是宁德的报表线条稀疏，
解析器把一列切成几列，表头和数字错开一格，期间列永远认不出来。我们为此写了稀疏列合并、续表链接和折行合并三项处理，佐证率升到 93.0\%，错配始终为 0（第 4 章）。
教训是：虚构样例只能做回归测试，不能代表真实数据上的效果。
\end{story}

\begin{story}[40 次会话里有 9 次没有正常结束]
五个工作区共 40 次会话，9 次退出码非 0，占 22.5\%，这些会话的费用合计 39.94 美元，占总费用 169.71 美元的 23.5\%。最贵的一次是宇树科技的排版会话，
跑了 77.6 分钟、花了 18.07 美元，最后没有正常结束。好在编排器按闸门状态判断进度，不看会话退出码：下一次会话从未通过的闸门继续，已经写出的文件保留，
所以这些会话的工作没有全部浪费。但这确实是成本波动的主要来源，也是 1005 批次比早期工作区便宜的主要原因。
\end{story}

\begin{story}[护栏误报，模型就去改写正确的句子]
护栏曾把“茅台酒销售公司净利率 33.35\%”当成公司整体净利率核对，把“第二季度收入”当成全年收入核对。主笔会话的应对是改写句子绕开护栏。
句子是改过去了，但信息变少了。我们后来让护栏识别分部、子公司和单季度的表述（第 7 章）。教训是：护栏的误报和漏报一样需要修，只是代价更隐蔽。
\end{story}

\begin{story}[评审嫌“不是……而是……”太多]
人工质检时，复核人发现评审把“出现几次‘不是……而是……’”当作机器腔，判了不通过，认为“过于严格”。这个句式确实是模型写作的常见痕迹，
我们自己也在复核编辑的技能说明里把它列为要改的句式；但评审的标准应该是“重复、无意义”，而不是“出现过”。评分标准 v1 据此改了措辞（第 9 章）。
\end{story}

\begin{story}[程序评级“卖出”，22 位分析师全部看多]
招商轮船 1005 版研报的程序评级是“卖出”，DCF 基准每股价值 9.55 元，较现价低 52.6\%，而资料包里 22 位卖方分析师全部给出买入或强推。
主笔按程序评级写，把分歧落在 2027 年以后的运价与新船交付；复核清单第 2 条请人核对 2027～2030 年的收入与毛利率假设是否过于悲观。
我们没有改评级，因为程序评级只是假设的函数；但它提醒我们，以五年期现金流为主的 DCF 对强周期公司未必是合适的锚，P/B 或重置成本估值尚未实现。
\end{story}

\begin{story}[删掉一节，评审说“只有标题没有内容”，却仍判通过]
评审校准时，我们删掉两份研报的“跟踪指标与结论失效条件”一节，评审在理由里指出这一节是空白页，却因为全文别处散落着跟踪指标而判了通过。
这件事有两层含义：劣化算子没把内容删干净（图表里的同类内容由排版脚本生成，不在文稿里），评审条目对结构完整性也不敏感。
改进劣化算子列在第 13 章的待做事项里。
\end{story}

\section{人工参与的方式}

系统并不排斥人。人在三个位置介入：给定输入（股价及来源、可比公司清单）；阅读人工复核清单，处理系统拿不准但重要的事项；
做质检复核，判断评审意见是否成立。人工质检用的是一个导出的静态网页，列出评审判“不通过”的条目、理由和证据，复核人逐条选择“确认”“误报”或“待查”并写批注，
结果回收到 \file{eval/calibration/human/}，用于计算评审精确率和修订评分标准。
